\section*{Executive Summary}

The dark-matter paradigm explains galactic rotation curve anomalies by postulating a new, non-luminous matter component. Despite decades of null results in direct, indirect, and collider searches, the hypothesis persists largely because no comparably rigorous alternative has been tested against the full breadth of observational data. This paper develops and tests such an alternative.

We keep the Einstein tensor standard (general relativity) and do not assume particulate cold dark matter (CDM). Instead, we use observed galactic phenomenology---rotation curves from the SPARC database of 175 galaxies---to infer the effective potentials and effective stress--energy required on the right-hand side of the Einstein equation, then compress that freedom with a low-degree-of-freedom ``intrinsic medium response'' closure.

The organizing principle is to treat the Einstein equation as an inverse problem,
\begin{equation}
G_{\mu\nu} = 8\pi G\left(T^{\mathrm{bar}}_{\mu\nu} + T^{\mathrm{resp}}_{\mu\nu}\right),
\end{equation}
where $T^{\mathrm{resp}}_{\mu\nu}$ is the response-sector stress--energy required by the observed phenomenology and then constrained by a low-degree-of-freedom closure.

\begin{itemize}
\item explains galaxy halo phenomenology without particulate CDM, and
\item admits a cosmological completion consistent with the empirical $\Lambda$CDM success set.
\end{itemize}

The response sector is activated near a baryonic boundary layer, produces an asymptotic tail in the extra acceleration, and is described by $\le 2$ degrees of freedom per galaxy. Across 175 SPARC galaxies, 92\% show very strong improvement over baryons-only, 97\% pass BIC overfitting tests, and the baryonic Tully--Fisher relation emerges as a consequence rather than an imposed prior.

We present explicit falsification criteria, discriminant tests (including lensing-facing statements and cosmological viability requirements), and a domain-partitioned roadmap from galaxy phenomenology through gravitational slip to large-scale structure. Claims are partitioned by domain: this work establishes the galaxy-domain closure and enumerates the cosmological completion requirements and discriminants.
